% ------------------------------------------------------------------------------

\begin{figure}[htpb]
    \lstinputlisting[language=csh,frame=single,basicstyle=\footnotesize\ttfamily]{fluent.sh}
    \caption{Source code for \file{fluent.sh}}
	\label{fig:fluent.sh}
\end{figure}

The job script in \xf{fig:fluent.sh} runs Fluent in parallel over 32 cores. 
%Of note, we have requested e-mail notifications (\texttt{-m}), are defining the 
Of note, we have requested e-mail notifications (\texttt{--mail-type}), are defining the 
%parallel environment for, \tool{fluent}, with, \texttt{-sgepe smp} (\textbf{very 
parallel environment for, \tool{fluent}, with, \texttt{-t\$SLURM\_NTASKS} and \texttt{-g-cnf=\$FLUENTNODES} (\textbf{very 
important}), and are setting \api{\$TMPDIR} as the in-job location for the
``moment'' \file{rfile.out} file (in-job, because the last line of the script 
copies everything from \api{\$TMPDIR} to a directory in the user's NFS-mounted home). 
Job progress can be monitored by examining the standard-out file (e.g.,
%\file{flu10000.o249}), and/or by examining the ``moment'' file in 
\texttt{slurm-249.out}), and/or by examining the ``moment'' file in 
\texttt{/disk/nobackup/<yourjob>} (hint: it starts with your job-ID) on the node running
the job. \textbf{Caveat:} take care with journal-file file paths.


